% ===== 10_discussion =====
\section{讨论与结论}
\label{sec:discussion}
\subsection{主要发现与相关研究}
同一社区非隔离感染人数上限下，接触减少与追踪隔离的分配会改变控制时长、累计新增感染和未感染者隔离数量。维持 $I=\eta$ 只约束乘积 $c(1-q)$，而追踪隔离还将未感染接触者移出社区。在固定二维进入状态 $(S^*,\eta)$ 和退出状态 $(S_c,\eta)$ 下，较大的 $c_c(S)$ 配合较高的隔离比例，使 $S$ 下降更快。由定理\ref{thm:joint:comparison}，当 $0<\beta<1$ 时，控制期累计新增感染随控制时长增加，新增隔离易感者人数则减少。因此，无额外能力限制时，仅隔离控制最快且控制期累计新增感染最少，代价是最多未感染者进入隔离。

能力界决定能否完成这一过程，以及每个状态上的允许分配区间。整段可行时，最大允许接触率 $\overline c(S)$ 给出最短控制时长，并在 $0<\beta<1$ 时给出最少的控制期累计新增感染。二次成本则在允许区间内同时权衡干预强度与时间消耗，权重改变其中的成本选择。命题\ref{prop:joint:terminal-allocation}表明，无额外能力限制，或整段可行且 $c_{\min}<c_0$、$q_{\rm cap}>q_0$ 时，最小成本分配在退出附近共同使用两种措施。与某一种单控制的严格成本比较还要求该单控制整段可行。

表\ref{tab:joint:compare}的两组同条件比较显示，最小成本分配相对仅隔离控制降低了 $J$，同时延长控制期并增加累计新增感染。共同完整初值、首次触发规则、阈值 $\eta>1$ 和前后常规控制使控制期感染排序延伸到各自的清零终点，见推论\ref{cor:joint:whole-epidemic}。有限能力代表点和权重网格展示了允许区间与成本选择的不同作用，数值切换仍由全部候选的成本比较确定。局部端点判据、有限权重网格和离散乘子点分别保留第\ref{sec:joint:cost}、\ref{sec:joint:weights-numerics}、\ref{sec:joint:duration-numerics}节的适用范围。

已有研究分别提供了恒定感染策略、分段干预解析和数据重构参照。Miclo 等\cite{Miclo2022ICU}在 SIR 模型、ICU 上限和特定线性成本下得到包含恒定感染阶段的解析最优策略。Zhou 等\cite{Zhou2024SPCI}以分段常数接触率和隔离比例研究干预时机，并用首次积分连接不同阶段。He 等\cite{He2023TDINN}从疫情数据重构随时间变化的接触率和隔离比例，为西安情景提供外部参照。本文在已定义的阈值控制类中，以状态参数化比较两种措施的全部允许分配，分别评价感染流入、未感染者隔离和二次强度成本。

退出要求也决定可达到的感染结局。命题\ref{prop:exit-infection-lower-bound}给出从初始时刻到有限时刻达到 $S=S_c$ 所需的累计新增感染下界，重新分配两种干预不能消除这一限制。对于仅隔离控制，降低阈值使感染峰值更低，却延长维持阈值直至同一易感者水平 $S=S_c$ 的过程。西安外部参照中，仅隔离控制的二次成本低于 TDINN 重构，感染峰值、累计新增感染和清零时间均较高，见表\ref{tab:xian_summary}。两者的阈值与退出要求不同，因而较低的 $J$ 只表示所定义归一化成本较低，不能代替对感染规模和终止任务的评价。

在当前无回流模型中，阈值控制于 $S=S_c$ 恢复常规，随后社区感染人数继续下降。按 $I=1$ 停止积分则给出操作性清零终点，恢复常规后的增长方向仍由该时刻的易感者状态决定。在西安 TDINN 重构的这一终点，若直接恢复常规接触率和隔离比例，状态回代给出 $\beta c_0(1-q_0)S/(\gamma N)\approx4.42>1$。该值是名义参数下的终点状态检查，未模拟后续反弹轨迹。本文保留其原有清零参照，并将其与达到规定退出状态的任务分别解释。

\subsection{适用范围}
\label{sec:dom:limits}
有效人口影响达到退出状态所需减少的易感者规模，但第\ref{sec:dominance}节的人口情景沿用全市拟合的传播参数与绝对初值，未对不同人口重新拟合。接近人口必要下界 $N_{\rm floor}$ 时，$S/N$ 对 $1$ 的偏离增大，原有 $S/N\approx1$ 条件下估计的参数能否移植需要另行判断。实际干预还可能改变混合范围，阈值控制与 TDINN 未必对应同一个有效人口。真实 $N_{\rm eff}$ 需由接触调查、出行信息或防控单元等外部资料确定，相对于固定全市参照的人口情景只能作条件性解释。

初值的缩放方式还影响启动和清零时间。既有近似 $t_1\approx r^{-1}\ln(\theta/i_0)$，其中 $r=\beta c_0(1-q_0)-\gamma$，说明小初值对启动时刻的影响。固定绝对初值时，归一化初值随人口改变；固定归一化初值的规模结论不能替代这一比较。清零终点采用绝对人数 $I=1$，因此时间指标仍含人口尺度。已有初值设定比较中，清零时间边界比成本、控制时长和累计新增感染边界更敏感。在所考察的两种初值设定及数值范围内，未找到同时满足人口相容下界与清零时间参照的参数点；连续允许人口集合仍未确定。完整人口比较的条件与有限证据见补充材料第\ref{sec:sup:population}节。

\FloatBarrier

追踪隔离缩短控制期的机制依赖于隔离易感者不返回社区。允许 $S_q$ 回到 $S$ 后，社区易感者可能不再单调下降，恢复常规后的传播还受回流影响。因此，以 $S$ 参数化控制、由时长确定感染排序及在 $S=S_c$ 退出的后续分析，都需在回流模型中重新建立。西安低阈值情景的长时结果沿用固定参数和无回流假设，适合作为解析基准，其时间估计不能直接解释为实际长期政策效果。

预先生成的时间开环控制依据名义模型维持 $I=\eta$，数值积分验证的是同一模型下的实现。观测误差会影响首次触发的判断，参数失配或执行延迟会使实际感染流入与移出偏离平衡，开环函数本身没有对偏离进行纠正的机制。本文未建立这些误差下的稳定性或峰值保证，提前干预与安全余量需另行评价。监测资料还需估计非隔离感染人数这一存量，某日新增确诊数不能直接替代 $I$。当前能力界限制控制强度，未限制调整速度。

二次成本\eqref{eq:cost:J}计量相对于常规水平的额外干预强度，其权重没有换算为货币或隔离资源价格。隔离比例恢复 $q_0$ 后，这一额外强度成本为零，既有隔离人群却可能继续占用资源。实际负担还取决于隔离人数、隔离时长和计费方式。若在目标中计入隔离人天，或加入累计预算、控制变化率限制，不同状态上的选择可能相互影响，式\eqref{eq:joint:minimum-cost}的逐状态分离需重新分析。边界明确的校园、厂区与依靠分区形成的传播单元也有不同成本，本文未计入建立和维持分区的投入。$J_c=0$ 仅表示模型中的接触率未进一步降低。

社区感染上限 $\eta$ 是外生设计参数。将它用于当地医疗容量管理，需结合 $I$ 与 $I_q$ 两类感染者的重症风险、入院时间和住院过程，建立感染流入与资源占用的联系。补充材料第\ref{sec:sup:medical-link}节的充分条件采用共同医疗需求等附加假设，第\ref{sec:sup:threshold-scale}节只提供阈值量级参考；二者均不构成西安当地的临床标定。对 $I$ 的限制也不能直接作为对 $I+I_q$ 或 ICU 占用的同一上限。本文数值结论沿用既有参数组、有限能力代表点和有限权重范围，其他阈值、人口或退出任务下的分配仍需分别评价。

\subsection{结论}
在无隔离易感者回流的 SIQR 型模型中，首次触发阈值、维持 $I=\eta$ 和在 $S=S_c$ 恢复常规控制构成本文的阈值控制任务。仅隔离控制给出启动、控制轨迹、退出后动态及累计新增感染的解析基准。在触发区间内，提高阈值推迟启动，降低最大隔离比例、控制时长和成本；当 $\eta>1$ 时，计算至清零终点的累计新增感染随之增加。联合控制由接触率状态函数 $c_c(S)$ 表示，能力界给出整段可行的充要条件。固定二维进入与退出状态且 $0<\beta<1$ 时，控制期累计新增感染与控制时长同序，新增隔离易感者人数反序，最大允许接触率配合隔离达到最快和感染最少的分配。

对于本文的二次成本及不含额外跨状态约束的允许控制类，最小值由逐状态比较端点与有限驻点达到。在触发且整段可行的阈值范围内，最小成本随阈值严格下降，能力扩大则使其不增。无额外能力限制时，最小成本分配在退出附近共同使用两种措施，其成本严格低于两种单控制各自的成本。规定有限时刻达到 $S=S_c$ 还产生累计新增感染下界。因而，同一峰值目标下应分别评价控制时长、累计新增感染、隔离负担与二次成本，退出任务也是这些评价条件的一部分。


\FloatBarrier

